\input{../tex-inputs/latex-korrekturansicht-vorspann}

\section[Arthur Schnitzler an Theodor Herzl, 11. 5. 1893]{L03943 Arthur Schnitzler an Theodor Herzl, 11. 5. 1893}
\nopagebreak\mylabel{L03943v}
\rehead{ }\normalsize\beginnumbering\briefempfaengerindex{Herzl, Theodor@\textsc{Herzl, Theodor}!zzzSchnitzler, Arthur@\emph{von Arthur Schnitzler}!1893-05-111@{11. 5. 1893}|(be}
\toendnotes[C]{\smallbreak\pagebreak[2]}
\correspDesc{Versand  durch Arthur Schnitzler am 11. 5. 1893 in Wien
\newline{}Erhalt  durch Theodor Herzl im Zeitraum [12. 5. 1893
                  – 16. 5. 1893?] in Paris}\toendnotes[C]{\smallbreak}
\Standort{Jerusalem, Central Zionist Archives, H1:1924-5.}
\physDesc{Brief, 1 Blatt, 3 Seiten, 597 Zeichen (Briefpapier mit Trauerrand)
\newline{}Handschrift: schwarze Tinte, deutsche Kurrent
\newline{}Ordnung: mit Bleistift von unbekannter Hand innerhalb das Konvoluts nummeriert: »(5)« und paginiert:
                                    »19« }\toendnotes[C]{\smallbreak}
\pstart{}{\pb}Verehrteſter Freund,\pend\vspace{0.5em}
\pstart
           innigen Dank für Ihre freundliche \textcolor{blue}{Teilnahme}\pwindex{Schnitzler, Johann 10.\,4.\,1835 Nagykanizsa – 2.\,5.\,1893 Wien@\textsc{Schnitzler, Johann} (10.\,4.\,1835 Nagykanizsa – 2.\,5.\,1893 Wien), \emph{Laryngologe}|pwv}{}\ledrightnote{{$\rightarrow$}\emph{\textcolor{blue}{Johann Schnitzler}}}; \substVorne{}\textsuperscript{d}\substDazwischen{}m\substHinten{}eine Angehörigen danken Ihnen gleichfalls aufs wärmſte! Man muſs ja auf alles
               mit den gleichen Worten danken, aber gerade ſo wie bei den Ausdrücken Ihres
               Mitgefühls manche Saiten in mir {\pb}mitvibrirten, darf ich wohl
               auch hoffen, daſs Sie aus den Worten meines Dankes mehr herausleſen, als das kühle
                  \label{K_L03943-1v}\edtext{\textsc{p. r.}}{\lemma{\textnormal{\emph{p. r.}}}\Cendnote{\textnormal{pro recipiendo (lateinisch): für den
                  Empfang}}}\label{K_L03943-1}, das für alle ist!\pend
           
\pstart
           Haben Sie Ihr Verſprechen ſchon ganz vergeſſen? Ihr heutiges \label{K_L03943-2v}\edtext{\textcolor{green}{\textsc{Feu{[}{]}lleton}}\pwindex{Herzl, Theodor 2.\,5.\,1860 Budapest – 3.\,7.\,1904 Edlach@\textsc{Herzl, Theodor} (2.\,5.\,1860 Budapest – 3.\,7.\,1904 Edlach), \emph{Schriftsteller, Journalist}!Firniß. Bilder aus dem Pariser Kunstleben. Karl VII. in Chinon@\strich\emph{Firniß. Bilder aus dem Pariser Kunstleben. Karl VII. in Chinon}|pwv}{}\ledrightnote{{$\rightarrow$}\emph{\textcolor{green}{Firniß. Bilder aus dem Pariser Kunstleben. Karl VII. in Chinon}}}}{\lemma{\textnormal{\emph{Feulleton}}}\Cendnote{\textnormal{\textcolor{blue}{Theodor Herzl}\pwindex{Herzl, Theodor 2.\,5.\,1860 Budapest – 3.\,7.\,1904 Edlach@\textsc{Herzl, Theodor} (2.\,5.\,1860 Budapest – 3.\,7.\,1904 Edlach), \emph{Schriftsteller, Journalist}|pwk}: \emph{\textcolor{green}{Firniß. Bilder aus dem Pariser Kunstleben. Karl VII. in
                        Chinon}\pwindex{Herzl, Theodor 2.\,5.\,1860 Budapest – 3.\,7.\,1904 Edlach@\textsc{Herzl, Theodor} (2.\,5.\,1860 Budapest – 3.\,7.\,1904 Edlach), \emph{Schriftsteller, Journalist}!Firniß. Bilder aus dem Pariser Kunstleben. Karl VII. in Chinon@\strich\emph{Firniß. Bilder aus dem Pariser Kunstleben. Karl VII. in Chinon}|pwk}}. In: \emph{\textcolor{green}{Neue Freie Presse}\pwindex{Neue Freie Presse@\emph{Neue Freie Presse}|pwk}},
                     Nr. 10.314, 11. 5. 1893, Morgenblatt, S. 1–4. }}}\label{K_L03943-2} hat
               mir wieder gezeigt, {\pb}einen wie ſchönen Dienſt Sie mir eben
               jetzt mit den alten Stücken erwieſen!\pend
           
\pstart
           Herzlich der Ihre{\\[\baselineskip]}\spacefill\mbox{ArthSchnitzler}\pend
           \leftskip=0em{}
\pstart
           11. 5. 93{ }\textcolor{pink}{Wien}\oindex{Wien@\textbf{Wien}, \emph{Verwaltungsgebiet}|pw}{}\ledrightnote{\textcolor{pink}{Wien}}.\pend
           \selectlanguage{ngerman}\endnumbering\briefempfaengerindex{Herzl, Theodor@\textsc{Herzl, Theodor}!zzzSchnitzler, Arthur@\emph{von Arthur Schnitzler}!1893-05-111@{11. 5. 1893}|)be}\mylabel{L03943h}  \newcommand{\dateiname}{L03943}\newcommand{\titel}{Arthur Schnitzler an Theodor Herzl, 11. 5. 1893}\newcommand{\editorInnen}{Herausgegeben von Jahnke, SelmaMüller, Martin Anton}\input{../tex-inputs/latex-korrekturansicht-abspann}
